% GENERATED FILE. Edit canonical lessons, then run npm run generate:books.
% canonical-source-hash: fnv1a64:8c7ff3fb6f448f06
% canonical-lessons: TA-C38-udambu, TA-C38-sugam, TA-C38-vidai

\chapter{Keeping Well, and Leaving}
\label{ch:ta-keeping-well-and-leaving}
\hlchaptermodality{38}

\begin{chapteropening}
\textbf{By the end of this chapter:} \emph{I can ask \ta{உங்கள்} \ta{உடம்பு} \ta{எப்படி} \ta{இருக்கு}?, recognise \ta{சுகமா}? and place it as the Sanskrit-derived option beside native \ta{நலம்}, say \ta{விடை} \ta{பெறுகிறேன்}, and show that \ta{விடு}, \ta{வீடு} and \ta{விடை} are one root while \ta{ஊர்} and \ta{நண்பன்} are not.}

Close the visit with \ta{நன்றி}. \ta{விடை} \ta{பெறுகிறேன்}., put \ta{விடு}, \ta{வீடு} and \ta{விடை} in one dictionary entry, refuse the second \ta{விடை} that means 'bull' and came from Sanskrit, and refuse to join \ta{ஊர்} and \ta{நண்பன்} to the pattern just because they are nearby.
\end{chapteropening}

\section[uḍambu]{\ta{உடம்பு} (uḍambu) — \textquotedblleft{}body,\textquotedblright{} and how Tamil asks after your health}

\label{lesson:TA-C38-udambu}



\begin{quote}
Introduce a friend respectfully and name the ending. (\emph{Ivar eṉ naṇbar} — the respectful \textbf{-\ta{அர்}}.) Next you would ask how they are.
\end{quote}

\subsection*{You'll want to know: \ta{உடம்பு}}
\begin{quote}
\textbf{\ta{உடம்பு}} — \emph{uḍambu} — \textbf{body}
\end{quote}

The whole of it, not a part. You have \textbf{\ta{தலை}} (\emph{talai}) and \textbf{\ta{கை}} (\emph{kai}) already; this is what they belong to, and the word inside the question Tamil speakers actually ask:

\begin{quote}
\textbf{\ta{உங்கள்} \ta{உடம்பு} \ta{எப்படி} \ta{இருக்கு}} — \emph{uṅgaḷ uḍambu eppaḍi irukku} — \textbf{how is your body}, that is, \textbf{how are you keeping}
\end{quote}

Every piece is yours: \textbf{\ta{உங்கள்}}, \textbf{\ta{எப்படி}}, and \textbf{\ta{இருக்கு}} from \textbf{\ta{இரு}}. What is new is that the health question hangs on a \textbf{noun} rather than on you.

Tamil has a second way to ask it, through \textbf{\ta{உணவு}} — the noun standing on \textbf{\ta{உண்}}, \textquotedblleft{}to eat.\textquotedblright{} Asking whether somebody has eaten is asking whether they are all right.

\subsection*{The letters in this word}
\textbf{\ta{உ}} + \textbf{\ta{ட}} + \textbf{\ta{ம்}} + \textbf{\ta{பு}}: the standing vowel \emph{u}, \textbf{\ta{ட}}, then \textbf{\ta{ம்}} — \textbf{\ta{ம}} shut by the puḷḷi you have used since \textbf{\ta{வணக்கம்}} — and \textbf{\ta{பு}}. \textbf{\ta{ம}} was one of the first letters you formed with a pen, alongside retroflex \textbf{\ta{ண}}.

\begin{sounds}
The letters say \emph{u-ṭa-m-pu}; your mouth says \emph{u\textbf{ḍ}a\textbf{mb}u}. \textbf{\ta{ட}} sits between vowels, so it softens — the \emph{paḍi} rule. \textbf{\ta{ப}} follows the nasal \textbf{\ta{ம்}}, so it voices — the \emph{naṉṟi} rule that also made \emph{naṇbaṉ}. Two rules, one short word, neither written down.
\end{sounds}

\begin{cousinweb}
Tamil has another word for the body, \textbf{\ta{உடல்}} (\emph{uḍal}), and the Dravidian dictionary does not treat them as rivals: \textbf{\ta{உடல்}}, \textbf{\ta{உடம்பு}} and \textbf{\ta{உடலம்}} sit in the \textbf{same entry}, with Kannada's \textbf{\ka{ಒಡಲು}} beside them. Usage has divided them: \textbf{\ta{உடல்}} leans literary, \textbf{\ta{உடம்பு}} is what you say aloud about someone unwell.
\end{cousinweb}

\begin{grammarlens}[title={Grammar Lens: where the person goes in a health question}]
Two shapes Tamil uses to ask about a state:

\noindent
\begin{tabularx}{\linewidth}{@{}>{\raggedright\arraybackslash}X>{\raggedright\arraybackslash}X@{}}
\toprule
\textbf{the question} & \textbf{where the person sits} \\
\midrule
\textbf{\ta{உங்கள்} \ta{உடம்பு} \ta{எப்படி} \ta{இருக்கு}} & as a \textbf{possessor} \\
\textbf{\ta{உனக்கு} \ta{எத்தனை} \ta{வயது} \ta{ஆகிறது}} & in the \textbf{dative} \\
\bottomrule
\end{tabularx}

The formal age question does not make you the subject at all — \textquotedblleft{}\textbf{to you}, how many years are becoming.\textquotedblright{} The body question makes \textbf{\ta{உடம்பு}} the subject.
\end{grammarlens}

\subsection*{Your turn}
Say these aloud:

\begin{itemize}
\raggedright
  \item the two softenings — \textquotedblleft{}uṭampu … uḍambu\textquotedblright{}
  \item the question — \textquotedblleft{}uṅgaḷ uḍambu eppaḍi irukku\textquotedblright{}
  \item the parts, the whole, and the two forms of one root — \textquotedblleft{}talai, kai … uḍambu … uḍal\textquotedblright{}
\end{itemize}

\begin{tcolorbox}[breakable,colback=teal!4,colframe=teal!35!black,title={Before you move on}]
Ask someone respectfully how they are keeping, through the body, and name the two sound rules inside \emph{uḍambu}. (\emph{Uṅgaḷ uḍambu eppaḍi irukku}; \textbf{\ta{ட}} softens \textbf{between vowels}, \textbf{\ta{ப}} voices \textbf{after a nasal}.) Are \textbf{\ta{உடல்}} and \textbf{\ta{உடம்பு}} different words, and where does the person sit in the formal age question? (\textbf{No}, same entry; \textbf{in the dative}.)
\end{tcolorbox}
\section[sugam]{\ta{சுகம்} (sugam) — \textquotedblleft{}ease,\textquotedblright{} the word Tamil did borrow}

\label{lesson:TA-C38-sugam}



\begin{quote}
Health through the body, and the other body-word in its entry. (\emph{Uṅgaḷ uḍambu eppaḍi irukku}; \emph{uḍal}.) A second answer came from abroad.
\end{quote}

\subsection*{You'll want to know: \ta{சுகம்}}
\begin{quote}
\textbf{\ta{சுகம்}} — \emph{sugam} — \textbf{wellbeing, ease, comfort}
\end{quote}

The \textbf{\ta{க}} is doing its old trick: the letter is \emph{ka}, but between vowels it comes out \textbf{g}. \emph{Su-\textbf{g}-am}, not \emph{su-ka-m}.

\begin{quote}
\textbf{\ta{சுகமா}} — \emph{sugamā} — \textbf{keeping well}
\end{quote}

It is \textbf{not} the ordinary Tamil \textquotedblleft{}how are you,\textquotedblright{} which is \textbf{\ta{நல்லா} \ta{இருக்கீங்களா}}. \textbf{\ta{சுகமா}} is warmer and more Sanskritic, and which you hear depends on somebody's \textbf{\ta{ஊர்}} — the native place whose name so often ends in \emph{-ūr}.

\subsection*{The letters in this word}
\textbf{\ta{சு}} + \textbf{\ta{க}} + \textbf{\ta{ம்}}: \textbf{\ta{ச}} with the \emph{u} sign, a bare \textbf{\ta{க}}, \textbf{\ta{ம}} closed by the puḷḷi. Nothing is doubled — hold that against \textbf{\ta{மன்னிக்கவும்}}, where \textbf{\ta{ன்னி}} has \textbf{\ta{ன}} dotted shut and written a second time, and \textbf{\ta{க்க}} does the same to \textbf{\ta{க}}.

\begin{grammarlens}[title={Grammar Lens: where the dative is required}]
\textbf{\ta{சுகமா}} has no dative in it, and a family of Tamil states demands one: \textbf{\ta{எனக்குத்} \ta{தமிழ்} \ta{தெரியும்}}, \textquotedblleft{}to me, Tamil is known.\textquotedblright{} Put a noun there and the shapes part:

\noindent
\begin{tabularx}{\linewidth}{@{}>{\raggedright\arraybackslash}X>{\raggedright\arraybackslash}X@{}}
\toprule
\textbf{the phrase} & \textbf{what happened} \\
\midrule
\textbf{\ta{நான்}} to \textbf{\ta{எனக்கு}} & the pronoun \textbf{changes its body} \\
\textbf{\ta{நண்பர்}} to \textbf{\ta{நண்பருக்கு}} & the noun just \textbf{takes -\ta{உக்கு}} \\
\bottomrule
\end{tabularx}

So \textbf{\ta{நண்பருக்குத்} \ta{தமிழ்} \ta{தெரியும்}}, and the formal age question does the same: \textbf{\ta{உனக்கு} \ta{எத்தனை} \ta{வயது} \ta{ஆகிறது}}. \textbf{\ta{சுகமா}} does not: knowing and ageing come \emph{to} a person; being well is what a person \textbf{is}.
\end{grammarlens}

\begin{cousinweb}
\textbf{\ta{சுகம்}} is Sanskrit \textbf{\dv{सुख}} (\emph{sukha}), \textquotedblleft{}ease\textquotedblright{} — the apology lesson in reverse. There, \textbf{\ta{மன்னிக்கவும்}} stood on Tamil's own \textbf{\ta{மன்னி}} while the neighbours built their \textquotedblleft{}sorry\textquotedblright{} on Sanskrit \textbf{kṣamā}. Here Tamil borrowed \emph{sukha} and \textbf{kept \ta{நலம்}} as the everyday word; Malayalam took the same loan, \textbf{\ml{സുഖം}}, as its ordinary greeting.

The old parse of \emph{sukha} is \emph{su-} \textquotedblleft{}good\textquotedblright{} plus \emph{kha} \textquotedblleft{}aperture\textquotedblright{} — a chariot's axle-hole. \textbf{Treat that as contested}: the standard Sanskrit dictionary hedges it with \textquotedblleft{}said to be,\textquotedblright{} and offers \emph{su-stha}.
\end{cousinweb}

\subsection*{Your turn}
Say these aloud:

\begin{itemize}
\raggedright
  \item the softened \ta{க} — \textquotedblleft{}sugam,\textquotedblright{} not \textquotedblleft{}sukam\textquotedblright{}
  \item the state that needs a dative — \textquotedblleft{}naṇbarukkut tamiḻ teriyum\textquotedblright{}
  \item the state that does not, and the word kept beside it — \textquotedblleft{}sugamā … nalam\textquotedblright{}
\end{itemize}

\begin{tcolorbox}[breakable,colback=teal!4,colframe=teal!35!black,title={Before you move on}]
Why is \textbf{\ta{க}} heard as \emph{g}, and what did the noun do that \textbf{\ta{நான்}} does not? (\textbf{Between vowels}; it \textbf{takes -\ta{உக்கு}}, while \textbf{\ta{நான்}} becomes \textbf{\ta{எனக்கு}}.) Which word did Tamil borrow, which did it keep, and how firm is the \textquotedblleft{}good axle-hole\textquotedblright{} story? (\textbf{\ta{சுகம்}}; \textbf{\ta{நலம்}}; \textbf{contested}.)
\end{tcolorbox}
\section[viḍai]{\ta{விடை} (viḍai) — \textquotedblleft{}leave,\textquotedblright{} and the root under the house}

\label{lesson:TA-C38-vidai}



\begin{quote}
\emph{Sugam} — which language it came from, and which Tamil word stayed beside it. (Sanskrit; \textbf{\ta{நலம்}}.) You came through the door, took the \textbf{\ta{நாற்காலி}} — four legs, counted out — were given tea, said \textbf{\ta{என்} \ta{தவறு}} when the wrong cup arrived, and introduced a friend. One thing is left.
\end{quote}

\subsection*{You'll want to know: \ta{விடை}}
\begin{quote}
\textbf{\ta{விடை}} — \emph{viḍai} — \textbf{leave; permission to depart}
\end{quote}

Not the act of going — you have \textbf{\ta{போய்} \ta{வருகிறேன்}} for that — but the \textbf{leave itself}, taken with \textbf{\ta{பெறு}} (\emph{peṟu}), \textquotedblleft{}to obtain\textquotedblright{}:

\begin{quote}
\textbf{\ta{நன்றி}. \ta{விடை} \ta{பெறுகிறேன்}.} — \emph{naṉṟi. viḍai peṟugiṟēṉ.} — \textbf{Thank you. I take my leave.}
\end{quote}

It is \textbf{formal and slightly literary}; leaving a friend's house you would say \textbf{\ta{போய்} \ta{வருகிறேன்}}.

The \textbf{\ta{தயவுசெய்து}} you learned for asking has a family across the Dravidian languages, all built on the same \emph{daya}, \textquotedblleft{}kindness,\textquotedblright{} plus a light verb. Leaving is not.

\subsection*{The letters in this word}
\textbf{\ta{வி}} + \textbf{\ta{டை}}. \textbf{\ta{வி}} is \textbf{\ta{வ}} with the \textbf{\ta{ி}} sign — the first vowel sign you learned, which made \textbf{\ta{நன்றி}} possible. \textbf{\ta{டை}} is \textbf{\ta{ட}} with the \emph{ai} sign \textbf{\ta{ை}}, which wraps around the front.

\textbf{\ta{நன்றி}} stands beside it above, so read it once more: \textbf{\ta{ந}} + \textbf{\ta{ன்}} + \textbf{\ta{றி}} — dental n, alveolar n dotted shut, hard \textbf{\ta{ற}} with that same \textbf{\ta{ி}}. And \textbf{\ta{பெறுகிறேன்}} has a \textbf{\ta{க}} between vowels, so it softens.

\begin{cousinweb}
\textbf{\ta{விடை}} is inherited Dravidian, and belongs to \textbf{\ta{விடு}} (\emph{viḍu}), \textquotedblleft{}to let go.\textquotedblright{} Leave is what you are \emph{let go} with.

Now the first word of these lessons. \textbf{\ta{வீடு}}, \textquotedblleft{}house,\textquotedblright{} is filed by the Dravidian dictionary in the \textbf{same entry as \ta{விடு}} — and \textbf{\ta{விடை}} is in it too:

\noindent
\begin{tabularx}{\linewidth}{@{}>{\raggedright\arraybackslash}X>{\raggedright\arraybackslash}X@{}}
\toprule
\textbf{the word} & \textbf{what it is} \\
\midrule
\textbf{\ta{விடு}} \emph{viḍu} & to let go, to release \\
\textbf{\ta{வீடு}} \emph{vīḍu} & the house — and, in the old sense, final release \\
\textbf{\ta{விடை}} \emph{viḍai} & leave to depart \\
\bottomrule
\end{tabularx}

You walked in through \textbf{\ta{வீடு}} and walk out through \textbf{\ta{விடை}}, one word in different clothes. Kannada's \textbf{\ka{ಬೀಡು}}, the same house with its \emph{b-}, sits on that root as well.

Two cautions, because a neat pattern is when to slow down. A second \textbf{\ta{விடை}} means \textquotedblleft{}bull,\textquotedblright{} a Sanskrit borrowing from \emph{vṛṣa}; and \textbf{\ta{ஊர்}} — the settlement inside Tanjore — and \textbf{\ta{நண்பன்}} are outside this root entirely.
\end{cousinweb}

\subsection*{Your turn}
Say these aloud:

\begin{itemize}
\raggedright
  \item the departure — \textquotedblleft{}naṉṟi. viḍai peṟugiṟēṉ.\textquotedblright{}
  \item the one root, three ways — \textquotedblleft{}viḍu … vīḍu … viḍai\textquotedblright{}
  \item the two that do not join it — \textquotedblleft{}ūr … naṇbaṉ\textquotedblright{}
\end{itemize}

\begin{tcolorbox}[breakable,colback=teal!4,colframe=teal!35!black,title={Before you move on}]
Take your leave formally, thanks in front of it, and name the earlier word that shares this root. (\emph{Naṉṟi. viḍai peṟugiṟēṉ.}; \textbf{\ta{வீடு}}, from \textbf{\ta{விடு}}, \textquotedblleft{}to let go.\textquotedblright{}) Do \textbf{\ta{ஊர்}} and \textbf{\ta{நண்பன்}} belong to that root? (\textbf{No} — a settlement and a closeness; not everything joins up.)
\end{tcolorbox}
